\documentclass[french, 11pt]{article}

\input{setup.tex}

\def\chapitre{9}
\def\pagetitle{Algorithmes pour l'intelligence artificielle.}

\begin{document}

\input{title.tex}

\renewcommand*{\E}{\mathbb{E}}
\renewcommand*{\P}{\mathbb{P}}

\begin{defi}{Apprentissage machine.}{}
    En apprentissage machine, on montre à la machine un grand jeu de données d'entraînement, et elle renvoie une fonction pour notre problème. On vérifie la qualité de cette fonction en la comparant à un nouveau jeu de données.
\end{defi}

\begin{defi}{Apprentissage supervisé.}{}
    Dans l'apprentissage supervisé, on utilise un jeu de données étiqueté, chaque élément est composé d'une intance, et d'une étiquette indiquant le résultat attendu.\\
    Un exemple : une image avec un texte la décrivant, ou un temps $t$ et une température. 
\end{defi}

\begin{defi}{}{}
    On distingue deux types de problèmes :
    \begin{itemize}
        \item Les problèmes de régression où l'ensemble des étiquettes est continu.
        \item Les problèmes de classification où l'ensemble des étiquettes est discret.
    \end{itemize}
    On se restreint aux problèmes de classification.
\end{defi}

\begin{defi}{Algorithme des $k$ plus proches voisins.}{}
    Étant donné une instance, on la compare à celles du jeu d'entraînement pour déterminer lesquelles lui ressembles le plus. On garde les $k$ plus proches. On renvoie alors l'étiquette la plus représentée parmi ceux-ci. La performance de l'algorithme dépend de la distance choisie pour calculer les plus proches voisins, et de $k$.\\
    Lorsque $k$ est petit, on ne considère que les voisins très proches, on est très sensible aux données d'entraînement.\\
    Il n'y a pas eu de généralisaiton, on parle de sur-apprentissage.\\
    Lorsque $k$ est trop grand, on considère des voisins qui ne sont plus représentatifs, on parle de sous-apprentissage.
\end{defi}

\begin{ex}{}{}
    Étant donné un nouvel élément, comment calculer ses $k$ plus proches voisins ?
    \tcblower
    \bf{Première solution.}\\
    Calculer la distance de notre nouvel élément à chacune des donénes, trier les données selon la distance, renvoyer les $k$ premiers, donc $O(n\log n)$.\\
    \bf{Seconde solution.}\\
    On stocke les $k$ meilleurs candidats, et pour chaque nouvel élément, on le compare au moins bon pour décider de le garder ou non.\\
    Avec un tableau on a une complexité en $O(nk)$. Avec une file de priorité, en $O(n\log k)$.\\
    En construisant un tas sur toutes les données : $O(n+k\log n)$.
    \bf{Troisième solution.} En $O(n+k)$...
\end{ex}

\begin{defi}{En $O(n+k)$.}{}
    On cherche à écrire un algorithme qui renvoie le $k^{\nt{ème}}$ plus petit élément d'un tableau. Une fois qu'on a cet élément, il suffit de parcourir toutes les données en ne conservant que les plus proches.\\
    \begin{algorithm}[H]
        \caption{Algorithme A.}
        \Entree{$T$ un tableau, $k$ le rang de l'élément recherché.}
        Couper le tableau en morceaux de taille 5.\\
        \PourCh{morceau}{Calculer sa médiane.}
        Construire $T_1$ rassemblant ces médianes.\\
        Calculer $x$ la médiane de $T_1$ avec $A(T_1,\frac{|T_1|}{2})$.\\
        Partitionner $T$ en deux tableaux $T'$ et $T''$, où $T''$ contient les éléments plus grands que $x$.\\
        Si $|T'|=k$, on renvoie $x$, si $|T'|>k$, on renvoie $A(T',k)$, sinon $A(T'', k-|T'|-1)$.\\
        Si $|T|<5$, on le trie et on renvoie le $k$-ème élément de $T$.
    \end{algorithm}
    \tcblower
    \bf{Terminaison.}\\
    Chaque appel récursif est réalisé sur un tableau strictement plus petit, les opérations qui ne sont pas des appels récursifs terminent.\\
    \bf{Correction.}\\
    Par récurrence sur $|T|$. Pour $|T|<5$ c'est bon.\\
    Supposons la correction sur tous les tableaux de taille inférieure à $n$. Soit $T$ un tableau de taille $n+1$.\\
    Chacun des appels récursifs est effectué sur un tableau strictement plus petit.
\end{defi}

\end{document}